\section{Noncommuting instruments with genuine feedback}\label{sec:quantum}
Let the physical Hilbert space have finite dimension $D_q\ge2$, and let $\mathcal H_0$ be its real space of traceless Hermitian matrices, with Hilbert--Schmidt norm. Put $d=D_q^2-1$. Let $\mu$ be uniform volume on $\{H\in\mathcal H_0:\norm H_{\HS}\le1\}$. Fix $0<\eta<1$, $0<\lambda<1$, and a finite family of unitaries $U_a$. At the start of each call the controller chooses $a$ from the past reports. The raw instrument is
\begin{equation}\label{eq:quantum-raw}
 \mathcal I_{a,H}(\rho)=T_a(K_H\rho K_H),\quad
 K_H=(I+\eta H)^{1/2},\quad
 T_a(\sigma)=\lambda U_a\sigma U_a^*
                   +(1-\lambda)\tr(\sigma)I/D_q.
\end{equation}
The weak measurement is followed by the chosen unitary and positive channel; it is not a free unitary chosen after seeing this call's report. Subsequent calls may react to that report. The report density is
$\ell(H,\rho)=1+\eta\tr(H\rho)$, which is genuinely state-dependent, and the effects generally do not commute.

The terminal task uses any fixed finite informationally complete POVM and its categorical Brier loss. Its probability vector is an injective affine chart of the density matrices. For instance, if $B_j$ is a Hilbert--Schmidt orthonormal traceless basis, then $\norm{B_j}_{\rm op}\le1$ and the $2d$ effects $(I\pm B_j)/(2d)$ are a POVM whose probabilities separate states. This is one randomized executable measurement, not a joint observation of incompatible observables. The classical controller has no uncharged quantum workspace.

\begin{lemma}[Positive-congruence inverse]\label{lem:congruence}
For every positive definite density matrix $\rho$, the map
\[
 E\longmapsto\frac{\sqrt E\rho\sqrt E}{\tr(E\rho)}
\]
after restricting to $E>0$, $\tr E=D_q$, is a smooth diffeomorphism onto the positive definite density matrices. Its inverse at $\sigma$ is
\begin{equation}\label{eq:inverse}
 B=\rho^{-1/2}(\rho^{1/2}\sigma\rho^{1/2})^{1/2}\rho^{-1/2},
 \qquad E=\frac{D_q B^2}{\tr B^2}.
\end{equation}
Its Jacobian and inverse Jacobian are bounded on compact positive subsets of these domains.
\end{lemma}
\begin{proof}
$B$ is the unique positive solution of $B\rho B=\sigma$, by taking the positive square root after conjugation by $\rho^{1/2}$. A normalized solution is therefore a positive scalar multiple of $B$, and the trace-$D_q$ constraint fixes the scalar in~\eqref{eq:inverse}. Matrix square root and inversion are smooth on the positive cone; the two displayed maps are smooth inverses. Compactness gives the derivative bounds and a positive minimum absolute determinant. Since $H$ is traceless, the effects $I+\eta H$ do lie in the required trace hyperplane. This is a finite-dimensional positive-matrix argument, not an unbounded-operator closure assertion.
\end{proof}

\begin{theorem}[Finite-state quantum feedback acquisition]\label{thm:quantum}
The instrument~\eqref{eq:quantum-raw}, with a fixed informationally complete terminal POVM, satisfies Theorem~\ref{thm:main} with $d=D_q^2-1$. It holds for every $0<\lambda,\eta<1$ and every fixed finite unitary action family. Thus the comparison is with the optimal full-history feedback choice of unitaries of the same experiment, not with a prescribed passive word. Constants may deteriorate near $\lambda=1$, $\lambda=0$, $\eta=0$, or loss of informational completeness.
\end{theorem}
\begin{proof}
Each raw map is completely positive and its integral is trace preserving: $\int H\,d\mu=0$ and $\tr T_a(\sigma)=\tr\sigma$. All finite future test probabilities are linear in the incoming density matrix; the terminal POVM separates it. Thus the predictive chart and unnormalized closure are obtained from the primitive instrument. Preparation of arbitrary density matrices in the sets below gives a Borel realization section. No spectral geometry is assumed as a resource law.

In density coordinates choose
\[
 D=\{\rho:\rho\ge(1-\lambda)I/(2D_q),\ \tr\rho=1\},\qquad
 K=\{\rho:\rho\ge(1-\lambda)I/D_q,\ \tr\rho=1\}.
\]
There is a positive buffer between $K$ and the relative boundary of $D$. Every update maps $D$ into $K$. The report density lies between $1-\eta$ and $1+\eta$, and
$\TV(\ell(\cdot,\rho),\ell(\cdot,\sigma))\le(\eta/2)\norm{\rho-\sigma}_{\HS}$.
The normalized map is jointly continuous. Lemma~\ref{lem:congruence}, the compact effect ball with $I+\eta H\ge(1-\eta)I$, and the invertible tangent map $\lambda U_a(\cdot)U_a^*$ imply a uniform bounded $d$-dimensional acquired density. Fixed affine POVM charts preserve these properties with their fixed norm and determinant constants.

It remains to derive feedback-uniform acquired overlap, rather than assume a transport coefficient. Write $\rho_* = I/D_q$. At $\rho_*$ the normalized weak measurement is exactly
$\rho_*+\eta H/D_q$. After $T_a$ it is $\rho_*+\lambda\eta U_aHU_a^*/D_q$. Hence its image contains the common ball
\[
 B_*:=B_{\HS}(\rho_*,\lambda\eta/(4D_q))
\]
with inverse effects of norm at most $1/4$. By the explicit inverse~\eqref{eq:inverse}, uniformly over this closed ball and the finite action family, there is $r_*>0$ such that every incoming $\rho$ with $\norm{\rho-\rho_*}_{\HS}\le r_*$ has an inverse effect in the ball of radius $1/2$ for every target in $B_*$. The Jacobian has a finite upper bound on this compact set and $\ell\ge1-\eta$. Thus every such acquired row dominates $q_*$ times uniform probability on $B_*$ for an explicitly defined positive $q_*$: take the minimum density given by~\eqref{eq:inverse} times $\operatorname{vol}(B_*)$, reduced to at most one if necessary.

Now restrict one report to $\norm H_{\HS}\le h$ with $\eta h\le1/2$. Functional calculus gives $\norm{K_H-I}_{\rm op}\le\eta h$. Since $\tr(K_H\rho K_H)\ge1/2$ and $\norm\rho_{\HS}\le1$,
\[
 \left\|\frac{K_H\rho K_H}{\tr(K_H\rho K_H)}-\rho\right\|_{\HS}
 \le8\eta h.
\]
For every chosen unitary the next distance to $\rho_*$ is therefore at most
$\lambda\norm{\rho-\rho_*}_{\HS}+8\lambda\eta h$.
Let $R=\sup_{\rho\in D}\norm{\rho-\rho_*}_{\HS}$. Choose $b\ge2$ and $h>0$ satisfying
\begin{equation}\label{eq:quantum-corridor}
 \lambda^{b-1}R\le r_*/2,\qquad
 \frac{8\lambda\eta h}{1-\lambda}\le r_*/2,\qquad
 h\le1,\quad\eta h\le1/2.
\end{equation}
Such choices exist for every fixed $\lambda<1$. Each quiet report has conditional probability at least $(1-\eta)h^d$. After $b-1$ such reports the state is in the $r_*$-ball, regardless of the adaptive actions. Proposition~\ref{prop:corridor} now gives
\[
 \epsilon_*=q_*\bigl((1-\eta)h^d\bigr)^{b-1}>0.
\]
Every quiet report and every unsuccessful corridor attempt remains physically counted. This derives the required overlap from actual kernels even when the improving channel is arbitrarily weak. Common report support and all remaining hypotheses were already verified. The theorem follows.
\end{proof}

\begin{proposition}[A strictly feedback-dependent qubit task]\label{prop:qubit-switch}
For $D_q=2$, the preceding theorem contains a fixed physical detector for which the optimal last unitary switches on a positive-mass acquired surface. Let $n_1,\ldots,n_4$ be tetrahedral unit Bloch vectors, let
\[
 a_1=\tfrac12,\quad a_2=a_3=a_4=\tfrac16,\quad t=\tfrac1{16},
 \qquad E_i=a_iI+t\,n_i\cdot\sigma.
\]
Use the terminal POVM $(E_i)$ and actions given by the identity and a $\pi$ rotation about an axis orthogonal to $n_1$. The two actions are strictly ordered with opposite signs according to $n_1\cdot r$, where $r$ is the current Bloch vector.
\end{proposition}
\begin{proof}
Each $E_i$ is strictly positive because $a_i>t$, and they sum to $I$. Since $\sum n_i=0$ and $\sum n_in_i^T=(4/3)I$, their outcome probabilities $p_i=a_i+t n_i\cdot r$ separate all three Bloch coordinates. Their Brier Bayes risk is
\begin{equation}\label{eq:qubit-risk}
 g(r)=\tfrac23-\tfrac{2t}{3}n_1\cdot r-\tfrac{4t^2}{3}\norm r^2.
\end{equation}
The weak-effect report ensemble is rotationally invariant. For a fixed eigenvalue radius $z\in[0,\eta/\sqrt2]$, integration over directions of the unnormalized L\"uders map gives mean Bloch multiplier $(1+2\sqrt{1-z^2})/3$. Radial averaging gives a constant $\gamma>1/3$ with $\E_\ell[r_{\rm post}]=\gamma r$. This follows directly by resolving a vector into the effect direction, whose multiplier is one, and its two transverse directions, whose multiplier is $\sqrt{1-z^2}$.

After the chosen unitary and depolarization, the linear term in~\eqref{eq:qubit-risk} is $-(2t/3)\lambda\gamma n_1\cdot R_ar$. The expected squared-norm term does not depend on the post-measurement rotation. The difference between identity and the flipping action is exactly
$-(4t/3)\lambda\gamma n_1\cdot r$.
The acquired row from the maximally mixed preparation has positive density on a ball about the origin, so both signs occur with positive probability. The control aligns a partially observed state with a fixed asymmetric detector; it is not a change of notation for an invariant task. The detector remains protected until commitment, and no classical controller is granted additional quantum memory.
\end{proof}
